% Generated draft: requires mathematical and source review.
\documentclass[11pt]{article}
\usepackage[margin=25mm]{geometry}
\usepackage{fontspec}
\setmainfont{Latin Modern Roman}
\usepackage{amsmath,amssymb,mathtools,unicode-math}
\setmathfont{Latin Modern Math}
\usepackage{enumitem,needspace,xurl,hyperref}
\setlength{\parindent}{0pt}
\setlength{\parskip}{5pt}
\setlength{\emergencystretch}{3em}
\newcommand{\N}{\mathbb{N}}
\newcommand{\Z}{\mathbb{Z}}
\newcommand{\Q}{\mathbb{Q}}
\newcommand{\R}{\mathbb{R}}
\newcommand{\C}{\mathbb{C}}
\newcommand{\F}{\mathbb{F}}
\newcommand{\A}{\mathbb{A}}
\newcommand{\PP}{\mathbb{P}}
\newcommand{\E}{\mathbb{E}}
\newcommand{\eps}{\varepsilon}
\newcommand{\dd}{\,\mathrm{d}}
\begin{document}
\section*{UnsolvedMath ID: 30000515}
\textbf{Source problem code:} catalogue entry
\subsection{Definitions and mathematical statement}
Use the usual complex numerical-Godeaux convention: $S$ is a smooth minimal projective surface of general type with $p_g(S)=h^0(S,K_S)=0$, $q(S)=h^1(S,\mathcal O_S)=0$ and $K_S^2=1$. The source calls it a Godeaux surface; this need not mean only the classical quotient construction.

Assume its bicanonical linear system $|2K_S|$ is a pencil with four distinct base points, and that every member is $2$-connected. The base points are points common to every divisor in the pencil; the precise scheme multiplicity/fixed-part conventions are not supplied. Here divisor $2$-connectedness means that for every decomposition of a member $C=A+B$ into nonzero effective divisors,
\[
 A\cdot B\geq2.
\]
It is not graph-theoretic or topological $2$-connectedness. The word \emph{every} includes special reducible or nonreduced members when they occur, not merely a general smooth member.

The proposed implication is
\[
 \bigl(|2K_S|\text{ has the stated four base points and all members are }2\text{-connected}\bigr)
 \quad\Longrightarrow\quad \pi_1(S(\C))=1.
\]
This uses topological simple connectivity in the complex setting; an \'etale-only variant would require separate clarification.

The supplied literature assessment checked 21 August 2026 says the reported construction's no-torsion conclusion is established, and that simply connected numerical Godeaux surfaces exist in other families, but does not verify simple connectivity of the particular surface in the cited construction. Absence of torsion (with the exact torsion invariant to be specified, usually Picard torsion here) is not itself triviality of the full fundamental group. Existence of other simply connected examples likewise does not establish this criterion. The cleaned question quantifies over surfaces satisfying its hypotheses, whereas the assessment discusses a particular construction; that scope must be checked against the source rather than silently reconciled.
\subsection{Short English statement}
Do the four-base-point bicanonical pencil and divisor-$2$-connectedness hypotheses force a numerical Godeaux surface to be simply connected? The supplied construction gives a no-torsion result, and other simply connected families are known, but neither alone verifies this surface or the general stated criterion.
\subsection{Sources}
\begin{itemize}
\item Local UnsolvedMath registry, permanent ID 30000515, exact source code catalogue entry; \emph{Simple Connectivity of Godeaux Surfaces}. Proposed year 2006, status partially\_solved; corrected\_verified statement specifies four distinct base points and every bicanonical member $2$-connected.
\item Supplied report citation: \emph{Classical Algebraic Geometry (2007)}, Oberwolfach Reports, DOI 10.4171/owr/2006/27, pp. 1615--1662, \url{https://doi.org/10.4171/owr/2006/27}. Title year differs from DOI/proposed year.
\item Supplied literature triage checked 2026-08-21: ``Schreyer--Stenger, simply connected Godeaux surfaces'', \url{https://arxiv.org/abs/2009.05357}; ``Reid, original construction'', \url{https://oa.tib.eu/renate/server/api/core/bitstreams/c73cc9ac-88d4-42ec-878f-d12ef7a6918c/content}; ``original Oberwolfach report'', \url{https://doi.org/10.4171/owr/2006/27}. This dated OWR-derived curation requests independent verification of time-sensitive status and cited claims.
\end{itemize}
\end{document}
